% Report del progetto S.A.G.E.: tengo qui struttura, sezioni e configurazioni.
% Per rigenerare il PDF basta compilare con pdflatex o latexmk -pdf.

\documentclass[11pt,a4paper]{article}
\usepackage[T1]{fontenc}
\usepackage[utf8]{inputenc}
\usepackage[italian]{babel}
\usepackage{lmodern}
\usepackage{microtype}
\usepackage{geometry}
\geometry{margin=2.5cm}
\usepackage{graphicx}
\usepackage[hidelinks]{hyperref}
\usepackage{booktabs}
\usepackage{array}
\usepackage{longtable}
\usepackage{float}
\usepackage{listings}
\usepackage{xcolor}
\usepackage{enumitem}

% Stile dei blocchi SQL nel report: colori leggibili e numeri di riga.
\definecolor{codepurple}{rgb}{0.58,0,0.82}
\definecolor{backcolour}{rgb}{0.95,0.95,0.92}
\definecolor{codegreen}{rgb}{0,0.6,0}
\definecolor{codeblue}{rgb}{0,0,0.8}

\lstset{
    language=SQL,
    backgroundcolor=\color{backcolour},
    commentstyle=\color{codegreen},
    keywordstyle=\color{codepurple}\bfseries,
    numberstyle=\tiny\color{gray},
    stringstyle=\color{orange},
    basicstyle=\ttfamily\small,
    breakatwhitespace=false,
    breaklines=true,
    captionpos=b,
    keepspaces=true,
    numbers=left,
    numbersep=5pt,
    showspaces=false,
    showstringspaces=false,
    showtabs=false,
    tabsize=2,
    morekeywords={LAST_INSERT_ID, ON, DUPLICATE, KEY, UPDATE, LIMIT, COALESCE, NULLIF}
}
\setlist{nosep, leftmargin=*}

% Dati del frontespizio, separati in comandi per modificarli in fretta.
\newcommand{\ProjectTitle}{S.A.G.E. - Sistema di Analisi e Gestione Economica}
\newcommand{\CourseName}{Corso: Basi di Dati}
\newcommand{\AcademicYear}{Anno accademico 2025/2026}
\newcommand{\GroupMembers}{Antonio Scharmuller \\ Riccardo Dallacasa}
\newcommand{\StudentNumbers}{Matricola 000000 \\ Matricola 000000}
\newcommand{\Emails}{antonio.scharmuller@studio.unibo.it \\ riccardo.dallacasa@studio.unibo.it}
\newcommand{\Supervisor}{Docente: Nome Cognome}
\newcommand{\UniversityLogo}{logo-unibo.png}

\begin{document}

\begin{titlepage}
    \centering
    \vspace*{1.5cm}
    {\Huge\bfseries \ProjectTitle \\[1.5em]}
    \vfill
    {\Large \CourseName \\}
    {\Large \AcademicYear \\}
    \vspace{1.5cm}
    {\large\bfseries Componenti del gruppo}\\
    \GroupMembers\\[1em]
    {\large\bfseries Matricole}\\
    \StudentNumbers\\[1em]
    {\large\bfseries Email universitarie}\\
    \Emails\\[1.5em]
    {\large \Supervisor}\\[2em]
    % Logo opzionale: basta mettere il file nella cartella o aggiornare il percorso.
    % \includegraphics[height=3cm]{\UniversityLogo}
    \vfill
    {\large \today}
\end{titlepage}

\pagenumbering{roman}
\tableofcontents
\newpage
\pagenumbering{arabic}

\section{Introduzione}
\subsection{Obiettivo del progetto}
La proposta iniziale di S.A.G.E. è la realizzazione di un sistema per la gestione e l'analisi delle finanze personali universitarie. L'obiettivo è supportare lo studente nella registrazione, organizzazione e analisi delle proprie transazioni economiche, offrendo strumenti per classificazione, pianificazione e monitoraggio nel tempo.

\subsection{Descrizione generale del sistema}
S.A.G.E. è progettato come un'applicazione multiutente in cui ogni studente dispone di un portafoglio finanziario personale isolato. Il sistema consente la registrazione di spese e entrate, la categorizzazione tramite categorie e tag, la definizione di budget temporali e la gestione di documenti digitali associati alle transazioni.

\subsection{Ambito del sistema}
L'ambito del progetto comprende:
\begin{itemize}
    \item gestione individuale del portafoglio economico personale;
    \item classificazione delle transazioni tramite categorie gerarchiche e tag trasversali;
    \item gestione dei budget mensili e per categoria;
    \item supporto per spese ricorrenti come modelli astratti;
    \item associazione di documenti digitali alle spese;
    \item analisi e reportistica aggregata.
\end{itemize}
Non sono previste funzionalità di condivisione di spese o conti comuni tra utenti.

\subsection{Utenti e ruoli previsti}
Il sistema prevede due ruoli principali:
\begin{itemize}
    \item \textbf{Utente}: gestisce le proprie transazioni, categorie e tag personali, definisce budget, registra documenti e visualizza report finanziari.
    \item \textbf{Amministratore}: gestisce gli elementi globali del sistema (categorie, tag, fonti) e analizza dati aggregati anonimi senza accedere ai dettagli finanziari personali degli utenti.
\end{itemize}

\subsection{Tecnologie utilizzate}
DBMS: MySQL\\
Linguaggio applicativo: Java\\
Interfaccia: Java Swing\\
Connessione DB: JDBC

\subsection{Scelte progettuali preliminari}
Fin dalla fase iniziale sono state adottate alcune scelte progettuali rilevanti: l'isolamento dei dati finanziari tra utenti, la distinzione tra elementi di classificazione globali e personali e la memorizzazione dei documenti digitali tramite percorso file anziché tramite BLOB nel database. Tali decisioni mirano a garantire privacy, personalizzazione e maggiore semplicità di gestione del DBMS.

\section{Analisi dei requisiti}
\subsection{Requisiti in linguaggio naturale}
Il sistema deve supportare la registrazione e il monitoraggio dei movimenti finanziari personali di uno studente universitario, distinguendo tra spese ed entrate. Ogni transazione deve avere importo, data e descrizione. Le spese devono essere associate a una categoria, mentre le entrate devono essere associate a una fonte. Il sistema deve permettere la definizione di budget mensili e per categoria, nonché la gestione di spese ricorrenti tramite un modello astratto. Deve inoltre conservare il percorso dei documenti associati alle spese senza memorizzare i file nel database.

\subsection{Analisi delle ambiguità e decisioni progettuali}
\begin{longtable}{p{3.5cm} p{5cm} p{6cm}}
\toprule
\textbf{Termine/Concetto} & \textbf{Criticità/Ambiguità} & \textbf{Decisione progettuale} \\
\midrule
\endfirsthead
\toprule
\textbf{Termine/Concetto} & \textbf{Criticità/Ambiguità} & \textbf{Decisione progettuale} \\
\midrule
\endhead
Periodo temporale & Come associare una spesa a un budget mensile? & Il periodo è derivato dalla data della transazione. L'entità \textsc{PERIODO} serve per la storicizzazione dei budget. \\
Spesa ricorrente & Generazione dei record nel DB? & È un modello astratto; la logica Java crea nuove \textsc{SPESE} a partire dai parametri del modello. \\
Privacy/Admin & L'admin vede i dati degli utenti? & Vincolo applicativo: l'admin accede solo a viste aggregate (COUNT/SUM), non ai dati granulari. \\
Documenti & Dove salvare i file? & Si salva solo il percorso sul file system; il DB non gestisce BLOB per evitare overhead. \\
\bottomrule
\end{longtable}

\subsection{Specifiche ristrutturate}
\begin{itemize}
    \item R1: Gestione delle transazioni finanziarie (spese/entrate) con dettagli minimi obbligatori.
    \item R2: Classificazione delle spese tramite categorie e classificazione trasversale delle transazioni tramite tag.
    \item R3: Definizione e storico dei budget mensili e dei budget per categoria.
    \item R4: Supporto per spese ricorrenti come modello di generazione automatica.
    \item R5: Salvare solo il percorso dei documenti digitali associati alle spese, senza memorizzare i file nel database.
    \item R6: Accesso amministrativo limitato a statistiche aggregate, senza visibilità sui dati finanziari personali.
\end{itemize}

\subsection{Glossario dei termini essenziali}
\begin{itemize}
    \item \textbf{Transazione}: entità astratta che rappresenta un movimento monetario con ID, data, importo e descrizione.
    \item \textbf{Budget}: reificazione della soglia di spesa legata a Utente, Categoria e Periodo.
    \item \textbf{Tag}: etichetta trasversale associabile a una o più transazioni, utile per analisi granulari non limitate alla sola categoria.
    \item \textbf{Categoria}: classificazione gerarchica 1:N applicata alle spese.
    \item \textbf{Fonte}: origine finanziaria obbligatoria per le entrate.
\end{itemize}

\subsection{Classi di utenza / viste del sistema}
\begin{itemize}
    \item \textbf{Utente/Studente}: accede al proprio portafoglio, registra transazioni, gestisce categorie personali e visualizza report.
    \item \textbf{Amministratore}: gestisce categorie/tag di sistema e visualizza statistiche aggregate.
\end{itemize}
Le viste del sistema includono il portafoglio personale, il cruscotto dei budget e la dashboard degli andamenti finanziari.

\subsection{Elenco delle operazioni principali}
\begin{itemize}
    \item Inserimento di una nuova spesa.
    \item Inserimento di una nuova entrata.
    \item Creazione/modifica/eliminazione di categorie personali.
    \item Creazione/modifica/eliminazione di tag personali.
    \item Definizione e aggiornamento dei budget mensili e per categoria.
    \item Gestione delle spese ricorrenti.
    \item Associazione di un documento digitale a una spesa.
    \item Visualizzazione di report aggregati e statistiche.
\end{itemize}

\subsection{Vincoli espliciti e impliciti}
\begin{itemize}
    \item L'amministratore può consultare solo dati aggregati (COUNT/SUM) e non i dettagli finanziari degli utenti.
    \item Ogni studente opera sempre sul proprio portafoglio personale: l'email dell'utente autenticato viene usata come chiave di sessione e come riferimento nelle tabelle dipendenti.
    \item I documenti vengono memorizzati solo tramite il percorso file, non come BLOB nel database.
    \item Ogni transazione deve contenere data, importo e descrizione.
    \item Il periodo del budget è derivato dal mese/anno della transazione e viene storicizzato tramite l'entità \textsc{PERIODO}.
    \item Il budget è opzionale: l'utente può registrare movimenti anche prima di aver definito una soglia di spesa.
    \item La spesa ricorrente è un modello astratto che genera transazioni effettive, non una transazione reale.
\end{itemize}

\section{Progettazione concettuale}
\subsection{Criteri generali di modellazione}

\subsection{Raffinamenti dello schema}
\subsubsection{Gerarchia delle transazioni (ISA)}
Abbiamo una specializzazione totale ed esclusiva tra spese ed entrate, tradotta nello schema relazionale con la relazione \textsc{TRANSIZIONE} e l'attributo \texttt{TipoTransazione}. Una spesa è associata obbligatoriamente a una categoria, mentre un'entrata è associata obbligatoriamente a una fonte. I tag sono collegati alle spese tramite la relazione \textsc{SPESA\_TAG}, perché servono a rendere più granulare la classificazione delle uscite.

\subsubsection{Modellazione ibrida (Sistema vs Personale)}
Per \textsc{CATEGORIA}, \textsc{TAG} e \textsc{FONTE} si applica una gerarchia ISA: elementi di sistema creati dall'admin e elementi personali creati dall'utente.
Nella traduzione logica questa scelta viene rappresentata con il flag \texttt{is\_system} e con l'attributo \texttt{Email\_Proprietario}. Le query applicative caricano gli elementi disponibili con la condizione \texttt{is\_system = TRUE OR Email\_Proprietario = ?}, in modo da mostrare insieme il catalogo comune e le personalizzazioni dello studente autenticato.

\subsubsection{Reificazione del budget}
Il budget è un'entità che mette in relazione utente e periodo, con un eventuale riferimento a una categoria. L’assenza della categoria rappresenta un budget mensile complessivo, mentre la presenza della categoria rappresenta un budget specifico. Questa progettazione consente di storicizzare i limiti di spesa e di fissare soglie per ambiti specifici. Nel capitolo di progettazione logica verrà valutata l'introduzione dell'attributo ridondante \texttt{Totale\_Speso\_Attuale} nel budget per migliorare le prestazioni di lettura nella dashboard.
Il sistema adotta una politica di budget opzionale: la registrazione di una spesa non richiede necessariamente l'esistenza preventiva di un budget. Se un budget è già presente, il suo \texttt{Totale\_Speso\_Attuale} viene aggiornato nella stessa transazione ACID dell'inserimento della spesa; se invece il budget viene creato in un secondo momento, il valore iniziale dell'attributo ridondante viene calcolato tramite una \texttt{SUM()} delle spese pregresse coerenti con periodo, utente ed eventuale categoria.

\subsection{Criticità rilevate}
\begin{itemize}
    \item La fonte deve essere legata esclusivamente a \textsc{ENTRATA}; una spesa ha una categoria di uscita.
    \item La spesa ricorrente punta a \textsc{CATEGORIA}, non a \textsc{SPESA}, perché è un modello generatore di tuple e non una transazione già avvenuta.
    \item L'utente viene identificato direttamente tramite \texttt{Email}, usata come chiave primaria e come chiave esterna nelle tabelle collegate. Questa scelta rende esplicito il legame tra sessione applicativa e isolamento dei dati.
\end{itemize}


\subsection{Schema scheletro E/R}

\subsection{Sviluppo delle viste o degli ambiti}
\subsubsection{Vista Studente e portafoglio personale}

\subsubsection{Vista classificazione delle transazioni}

\subsubsection{Vista budget e pianificazione}

\subsubsection{Vista spese ricorrenti e documenti}

\subsubsection{Vista amministratore e cataloghi globali}

\subsection{Raffinamenti proposti}

\subsection{Gerarchie presenti nello schema}

\subsection{Associazioni principali e cardinalità}

\subsection{Vincoli non rappresentabili direttamente nello schema E/R}

\subsection{Integrazione delle viste}

\subsection{Schema E/R finale}

\subsection{Dizionario di entità e associazioni}

\section{Specifiche funzionali}
\subsection{Elenco delle funzionalità richieste}

\subsection{Operazioni di inserimento}

\subsection{Operazioni di modifica}

\subsection{Operazioni di eliminazione}

\subsection{Operazioni di interrogazione}

\subsection{Operazioni statistiche / amministrative}

\subsection{Input, output e tabelle coinvolte per ogni operazione}

\section{Progettazione logica}
\subsection{Stima del volume dei dati}

\subsection{Descrizione delle operazioni principali e frequenza}

\subsection{Schemi di navigazione}

\subsection{Tabelle degli accessi}

\subsection{Analisi delle ridondanze}
L'attributo \texttt{Totale\_Speso\_Attuale} della relazione \textsc{BUDGET} è ridondante perché può essere ricavato sommando le spese presenti in \textsc{TRANSIZIONE}. La ridondanza viene comunque mantenuta per ridurre il costo delle letture più frequenti della dashboard, dove l'utente controlla spesso lo stato dei budget rispetto al limite impostato.

La coerenza viene garantita con due regole applicative:
\begin{itemize}
    \item quando viene inserita una nuova spesa, l'inserimento della transazione, l'associazione di tag/documento e l'aggiornamento degli eventuali budget esistenti avvengono nella stessa transazione;
    \item quando un budget viene creato o aggiornato dopo che alcune spese sono già state registrate, il valore iniziale di \texttt{Totale\_Speso\_Attuale} viene calcolato con una query di aggregazione sulle transazioni già presenti.
\end{itemize}

In questo modo il budget rimane opzionale senza perdere consistenza: l'assenza di un budget non blocca la registrazione della spesa, mentre l'esistenza di un budget impone la sincronizzazione del suo totale ridondante.

Per evitare che la ridondanza diventi una fonte di incoerenza, il progetto affianca al valore memorizzato anche la vista \texttt{v\_budget\_stato}, che calcola lo stato effettivo del budget tramite aggregazione sulle transazioni. Questa vista viene usata come riferimento logico per verificare il valore reale di spesa, residuo e superamento della soglia. Anche lo script di popolamento iniziale calcola \texttt{Totale\_Speso\_Attuale} con una \texttt{SUM()} sulle transazioni già inserite, in modo che i dati dimostrativi risultino coerenti anche se osservati direttamente da MySQL Workbench.

\subsection{Raffinamento dello schema}
\subsubsection{Eliminazione delle gerarchie}

\subsubsection{Soluzione logica anticipata}
La specializzazione totale ed esclusiva delle transazioni si traduce in una struttura logica collassata con la relazione \textsc{TRANSIZIONE}. La distinzione tra movimenti in uscita e movimenti in entrata è rappresentata da \texttt{TipoTransazione}, mentre i vincoli CHECK impongono la presenza della categoria solo per le spese e della fonte solo per le entrate.

\subsubsection{Eliminazione degli attributi composti}

\subsubsection{Gestione degli attributi multivalore}

\subsubsection{Scelta delle chiavi primarie}

\subsubsection{Traduzione delle associazioni}

\subsection{Traduzione di entità e associazioni in relazioni}
In questa fase si traducono le entità e le associazioni concettuali in relazioni logiche, specificando le chiavi primarie e le chiavi esterne necessarie per mantenere l'integrità referenziale.

\subsection{Schema relazionale finale}
Lo schema relazionale finale descrive le tabelle definitive, con i rispettivi attributi, tipi di dato e vincoli principali. Questa rappresentazione costituisce la base per la successiva fase di implementazione SQL.

\subsection{Vincoli di integrità}
In questa sezione si elencano i vincoli di integrità referenziali e di dominio che il database deve rispettare, come le chiavi primarie, le chiavi esterne, i vincoli UNIQUE e i controlli sui valori.

\subsection{Controllo della normalizzazione}
Qui si verifica che lo schema relazionale finale soddisfi i livelli di normalizzazione richiesti, identificando eventuali anomalie di inserimento, aggiornamento o cancellazione.

\section{Implementazione SQL}
\subsection{Script di creazione del database}

\subsection{Creazione delle tabelle}

\subsection{Definizione di chiavi primarie e chiavi esterne}

\subsection{Vincoli NOT NULL, UNIQUE, CHECK}

\subsection{Trigger, procedure o viste, se usati}

\subsection{Popolamento iniziale del database}
Il popolamento iniziale inserisce utenti demo, periodi, categorie, tag, fonti, transazioni, documenti, budget e spese ricorrenti. Per i budget dimostrativi, il valore ridondante \texttt{Totale\_Speso\_Attuale} non viene lasciato al valore di default, ma viene inizializzato tramite query di aggregazione sulle transazioni già presenti. In questo modo la base dati di prova è coerente sia nella visualizzazione applicativa sia nella consultazione diretta delle tabelle.

\subsection{Query SQL delle operazioni principali}
\subsubsection{Problemi e criticità}
Il progetto richiede di garantire l'integrità dei dati durante l'inserimento di una spesa o di un'entrata, e di gestire i budget in modo idempotente. Le principali criticità sono:
\begin{itemize}
    \item Isolamento utente: ogni query deve filtrare i dati tramite l'email dell'utente autenticato.
    \item Atomicità: l'inserimento della spesa, l'associazione dei tag, l'eventuale documento e l'aggiornamento dei budget devono avvenire nella stessa transazione.
    \item Divisione per zero: nelle statistiche percentuali l'assenza di entrate potrebbe causare errori se non si utilizza una protezione adeguata.
    \item Gestione del budget: la coppia chiave \texttt{(Email, ID\_Periodo, ID\_Categoria\_Key)} deve essere unica per permettere work-flow di upsert senza duplicati.
\end{itemize}

\subsubsection{Analisi}
Il sistema utilizza una traduzione collassata della gerarchia: la relazione \textsc{TRANSIZIONE} contiene sia spese sia entrate, distinte dall'attributo \texttt{TipoTransazione}. Il vincolo CHECK impone che una spesa abbia una categoria e nessuna fonte, mentre un'entrata abbia una fonte e nessuna categoria. Per i budget si introduce la logica di "upsert" per evitare duplicati e mantenere coerente l'attributo ridondante \texttt{Totale\_Speso\_Attuale}.

\subsubsection{Soluzione SQL}
\paragraph{U3 - Inserimento Atomico Spesa}
\begin{lstlisting}[caption={Inserimento di una spesa con tag e budget}]
START TRANSACTION;

INSERT INTO PERIODO (Mese, Anno)
VALUES (?, ?)
ON DUPLICATE KEY UPDATE ID_Periodo = LAST_INSERT_ID(ID_Periodo);

INSERT INTO TRANSIZIONE
(TipoTransazione, Importo, Data, Descrizione, Email, ID_Categoria, ID_Periodo, ID_Fonte)
VALUES ('S', ?, ?, ?, ?, ?, LAST_INSERT_ID(), NULL);

SET @id_spesa = LAST_INSERT_ID();

INSERT INTO SPESA_TAG (ID_Transizione, ID_Tag)
VALUES (@id_spesa, ?);

UPDATE BUDGET
SET Totale_Speso_Attuale = Totale_Speso_Attuale + ?
WHERE Email = ? AND ID_Periodo = ? AND ID_Categoria IS NULL;

UPDATE BUDGET
SET Totale_Speso_Attuale = Totale_Speso_Attuale + ?
WHERE Email = ? AND ID_Periodo = ? AND ID_Categoria = ?;

COMMIT;
\end{lstlisting}

\paragraph{U8 - Definizione/Aggiornamento Budget}
\begin{lstlisting}[caption={Upsert del budget con totale iniziale coerente}]
SELECT COALESCE(SUM(Importo), 0) AS Totale_Speso
FROM TRANSIZIONE
WHERE Email = ?
  AND ID_Periodo = ?
  AND TipoTransazione = 'S'
  AND (? IS NULL OR ID_Categoria = ?);

INSERT INTO BUDGET
(Importo_Limite, Alert_Soglia, Totale_Speso_Attuale, ID_Periodo, ID_Categoria, Email)
VALUES (?, ?, ?, ?, ?, ?)
ON DUPLICATE KEY UPDATE
Importo_Limite = VALUES(Importo_Limite),
Alert_Soglia = VALUES(Alert_Soglia),
Totale_Speso_Attuale = VALUES(Totale_Speso_Attuale);
\end{lstlisting}

\paragraph{A3 - Analisi Percentuale (protezione divisione per zero)}
\begin{lstlisting}[caption={Rapporto Spese/Entrate Mensile}]
SELECT
    (
        SUM(CASE WHEN TipoTransazione = 'S' THEN Importo ELSE 0 END)
        /
        NULLIF(SUM(CASE WHEN TipoTransazione = 'E' THEN Importo ELSE 0 END), 0)
    ) * 100 AS Percentuale_Spese_Su_Entrate
FROM TRANSIZIONE
WHERE Email = ? AND ID_Periodo = ?;
\end{lstlisting}

\subsubsection{Query di supporto e operazioni amministrative}
\paragraph{Cataloghi disponibili per utente}
\begin{lstlisting}[caption={Categorie, fonti e tag di sistema o personali}]
SELECT ID_Categoria, Nome, is_system
FROM CATEGORIA
WHERE is_system = TRUE OR Email_Proprietario = ?
ORDER BY is_system DESC, Nome;
\end{lstlisting}

\paragraph{Riepilogo delle transazioni per periodo}
\begin{lstlisting}[caption={Riepilogo transazioni per periodo}]
SELECT T.ID_Transizione, T.Importo, T.Data, T.Descrizione,
       COALESCE(C.Nome, F.Nome) AS Categoria_Fonte,
       CASE WHEN T.TipoTransazione = 'S' THEN 'Spesa' ELSE 'Entrata' END AS Tipo
FROM TRANSIZIONE T
LEFT JOIN CATEGORIA C ON T.ID_Categoria = C.ID_Categoria
LEFT JOIN FONTE F ON T.ID_Fonte = F.ID_Fonte
WHERE T.Email = ? AND T.Data BETWEEN ? AND ?
ORDER BY T.Data DESC;
\end{lstlisting}

\paragraph{Stato dei budget}
\begin{lstlisting}[caption={Budget con residuo e stato di superamento}]
SELECT ID_Budget, Mese, Anno, Ambito,
       Importo_Limite, Totale_Speso, Residuo, Superato
FROM v_budget_stato
WHERE Email = ?
ORDER BY Anno DESC, Mese DESC, Ambito;
\end{lstlisting}

\subsubsection{Motivazioni}
Atomicità applicativa: in Java, l'uso di \texttt{conn.setAutoCommit(false)} unito a queste query garantisce che spesa, tag, documento e aggiornamento budget vengano confermati o annullati insieme.\\
Flessibilità budget: il vincolo \texttt{UNIQUE(Email, ID\_Periodo, ID\_Categoria\_Key)} permette alla query \texttt{ON DUPLICATE KEY UPDATE} di aggiornare o creare il limite in modo efficiente.\\
Robustezza matematica: \texttt{NULLIF} trasforma lo zero in NULL, evitando l'errore di divisione per zero nel database engine.

\section{Progettazione dell’applicazione}
\subsection{Architettura generale dell’applicazione}
L'applicazione è organizzata secondo una struttura a livelli ispirata al modello MVC, integrata con il pattern DAO per l'accesso ai dati. Le classi della vista gestiscono l'interfaccia Swing, i controller raccolgono gli eventi e coordinano le richieste, i service contengono la logica applicativa e transazionale, mentre i DAO incapsulano le operazioni effettive sul database tramite JDBC.

Il flusso principale è quindi:
\[
\texttt{view} \rightarrow \texttt{controller} \rightarrow \texttt{service} \rightarrow \texttt{dao} \rightarrow \texttt{database}
\]
Questa separazione evita che l'interfaccia grafica dipenda direttamente dal database e rende più semplice modificare o testare le singole parti dell'applicazione.

\subsection{Tecnologie usate}
L'applicazione è sviluppata in Java e utilizza Swing per l'interfaccia grafica. La comunicazione con MySQL avviene tramite JDBC, con connessioni centralizzate nella classe di utilità dedicata. Gli script SQL di creazione, popolamento, viste, trigger e interrogazioni principali sono raccolti nella cartella \texttt{doc/sql}, separata dagli schemi concettuali e logici presenti in \texttt{doc/schemes}.

\subsection{Struttura dei package Java}
Il codice sorgente è collocato nel package \texttt{it.unibo.sage} ed è suddiviso nei seguenti sottopackage:
\begin{itemize}
    \item \texttt{model}: classi che rappresentano gli oggetti principali del dominio applicativo, come \texttt{Utente}, \texttt{Transazione}, \texttt{Budget}, \texttt{Categoria}, \texttt{Tag}, \texttt{Fonte} e \texttt{Documento};
    \item \texttt{dao}: interfacce DAO e implementazioni JDBC per separare le operazioni logiche sui dati dalla loro realizzazione sul DBMS;
    \item \texttt{service}: logica applicativa, gestione delle transazioni JDBC e coordinamento tra più DAO;
    \item \texttt{controller}: classi che collegano la vista ai service e mantengono stabile l'API usata dall'interfaccia;
    \item \texttt{view}: pannelli, frame e componenti grafici Java Swing;
    \item \texttt{utils}: classi di supporto, tra cui la gestione della connessione al database.
\end{itemize}

\subsection{Collegamento al database tramite JDBC}
La classe \texttt{DatabaseConnection} centralizza la creazione delle connessioni JDBC. I service aprono le connessioni necessarie e, quando una funzionalità coinvolge più tabelle, disattivano l'autocommit per garantire atomicità. Ad esempio, l'inserimento di una spesa, l'associazione dei tag, l'eventuale salvataggio del documento e l'aggiornamento dei budget avvengono nella stessa transazione applicativa.

Per la schermata dei budget, il service ricalcola lo stato di avanzamento partendo dalle transazioni effettive, così la vista applicativa resta allineata ai movimenti registrati anche quando un budget viene creato dopo l'inserimento di alcune spese.

\subsection{Organizzazione delle classi principali}
Le classi del package \texttt{model} rappresentano gli oggetti principali del dominio applicativo e vengono utilizzate per trasferire i dati tra DAO, service, controller e interfaccia grafica. La classe \texttt{Documento} rappresenta un file associato a una transazione, come ricevuta, fattura o prova di pagamento; nel database viene memorizzato il percorso del file e non il contenuto binario.

Per l'accesso ai dati è stato adottato il pattern DAO, separando le operazioni logiche sulle entità dalla loro implementazione concreta tramite JDBC. In questo modo i controller e i service non dipendono da query sparse nell'interfaccia, ma comunicano con classi dedicate alla persistenza.

\subsection{Gestione degli errori e validazione input}
Gli errori SQL vengono propagati verso i controller e gestiti dalla vista con messaggi all'utente. Le operazioni composte usano rollback in caso di eccezione, così da evitare stati parziali del database. Le validazioni di base, come campi obbligatori e importi non validi, sono gestite prima dell'invio dei dati ai service.

\subsection{Interfaccia grafica Java Swing}
L'interfaccia grafica è realizzata con Java Swing e organizzata in pannelli dedicati alle funzionalità principali: autenticazione, dashboard, movimenti, budget e navigazione laterale. I componenti grafici non eseguono query direttamente, ma invocano i controller, mantenendo separata la logica di presentazione dalla logica applicativa.

\subsection{Schermate principali}
Le schermate principali dell'applicazione sono raccolte nella cartella \texttt{doc/screenshots} e documentano i flussi più importanti: login, dashboard, registrazione dei movimenti, gestione dei budget e spese ricorrenti.

\subsection{Flussi d’uso principali}

Registrazione studente\\
Inserimento di una spesa\\
Inserimento di una entrata\\
Definizione di un budget\\
Gestione di una spesa ricorrente\\
Associazione di un documento digitale\\
Visualizzazione statistiche amministratore

\section{Test e validazione}
\subsection{Dati di test utilizzati}

\subsection{Test delle query SQL}

\subsection{Test delle funzionalità applicative}

\subsection{Casi limite gestiti}

\subsection{Errori noti o limitazioni}

\section{Conclusioni}
\subsection{Risultato ottenuto}

\subsection{Scelte progettuali principali}

\subsection{Possibili sviluppi futuri}

\appendix
\section{Script SQL completo}

\section{Query SQL complete}

\section{Screenshot dell’applicazione}

\section{Manuale utente sintetico}

\section{Eventuale link al repository}

\end{document}
